\scanpage{75}
\mspage{24}

$W$-asphériques. Pour de nombreuses \deleted{autres} \deleted{\unsure{équivalentes}} formulations équivalentes, je renvoie à P.S.

\uline{Proposition 2} Soit $X\mathrel{\overset{f}{\relbar}}Y\xrightarrow{g}Z$ dans $\mathrm{Cat}$. Si $f$ et $g$ sont \deleted{\unsure{des}} $W$-asphériques (\inserted{$W$-}coasphériques), de $=$ $gf$. Si $gf$ et $f$ sont $W$-asphériques ($W$-coasphériques), $g$ aussi.

\deleted{Rappel} \uline{Rappel} Un morphisme $p:S'\longrightarrow S$ de $\mathrm{Cat}$ est un \uline{isom. local} si $\forall\,s'\in\mathrm{Ob}\,S'$, le morphisme induit \[p_{s'}:S'_{/s'}\xrightarrow{\sim}S_{/s}\qquad(s=p(s'))\] est un isomorphisme : $p$ induit un isomorphisme entre les ``localisées'' de $S'$ et $S$ en $s'$ et $s=p(s')$ respectivement. Il est immédiat que cette notion est stable par composition, que si $gf$ et $g$ sont des isom locaux, alors $f$ l'est aussi, que tout isomorphisme est un isom. local. De plus, les isom. locaux sont \uline{stables par chgt de base}, ce qui résulte du fait \inserted{(trivial et important)} que si on a un carré cartésien de $\mathrm{Cat}$ \diagram{\begin{tikzcd}[ampersand replacement=\&] X\arrow[d,"f"']\& X'\arrow[l,"q"']\arrow[d,dash,"f'"]\\ S\& S'\arrow[l,"p"] \end{tikzcd}}, et si $x'\in\mathrm{Ob}\,X'$ a pour images $x,s',s$ dans $X,S',S$, alors le diagramme \diagram{\begin{tikzcd}[ampersand replacement=\&] X_{/x}\arrow[d,"{f_{/x}}"']\& X'_{/x'}\arrow[l,"{q_{/x'}}"']\arrow[d,"{f'_{/x'}}"]\\ S_{/s}\& S'_{/s'}\arrow[l,"{p_{/s'}}"] \end{tikzcd}} est également cartésien. Donc si $p_{/s'}$ est iso, $q_{/x'}$ aussi.

% archive sheet number: 75
